% ======================================================================
\section{Background and related literature}\label{sec:literature}

This section reviews the research this paper builds on. It covers how planetary boundaries are defined and translated to the national level, how social foundations are defined, what is known about the relationship between human well-being and emissions, the evidence on decoupling, and the choice between production- and consumption-based emissions accounting. It closes by identifying the gaps the paper addresses.

\subsection{Planetary boundaries and the safe operating space}

The planetary boundaries framework was introduced by \citet{rockstrom2009}. Its starting point is that the last ten thousand years, the Holocene, gave human societies an unusually stable climate and biosphere, and that human activity now threatens that stability. The framework identifies nine Earth-system processes that regulate the state of the planet: climate change, biosphere integrity, land-system change, freshwater use, biogeochemical flows of nitrogen and phosphorus, ocean acidification, atmospheric aerosol loading, stratospheric ozone depletion and the introduction of novel entities. For each it proposes a boundary, a level of human pressure beyond which the risk of large-scale, abrupt or irreversible change rises. The boundaries are set at the safe end of a ``zone of uncertainty''. Crossing a boundary does not trigger an immediate collapse, but it moves humanity into a zone of increasing risk.

\citet{steffen2015} updated the framework and introduced a two-level structure. Climate change and biosphere integrity are ``core'' boundaries, because each on its own could push the Earth system into a different state. The other seven operate largely through their effects on these two. \citet{richardson2023} provided the first assessment of all nine and concluded that six had been transgressed. For climate change, the boundary has two control variables: an atmospheric \COtwo{} concentration of 350 parts per million and an increase in radiative forcing of one watt per square metre. Both have been exceeded; the concentration boundary was crossed in the late 1980s \citep{fanning2022}.

Climate change is the focus of this paper for three reasons. It is a core boundary. Its main driver, \COtwo{} from fossil fuels and industry, is measured for almost every country every year with reasonable precision. And it can be translated into national terms through the idea of a carbon budget: the total amount of \COtwo{} that can still be emitted while keeping warming below a given level with a given probability \citep{ipcc2021spm}.

\subsection{From a global boundary to national shares}

Planetary boundaries are defined for the Earth as a whole. Comparing countries requires a rule for allocating the global safe operating space among them. This is not a purely scientific question, because any allocation embodies a principle of fairness. \citet{haeyhae2016} distinguish several such principles. One allocates according to population (equality). Others take account of historical responsibility for past pressures, of capability to reduce pressures, or of need.

The simplest and most widely used approach divides the global boundary equally among the world's population. \citet{oneill2018} applied this to the carbon budget. They took a budget of about 1{,}000 Gt \COtwo{} for 2011--2100, consistent with a 66\% chance of keeping warming below 2\textdegree C, and spread it evenly over the 90 years and across seven billion people. The result is an annual boundary of 1.61 t \COtwo{} per person. This paper uses that value as its main boundary. \citet{fanning2022} used a different, cumulative approach. They compared each country's emissions since 1850 with its population share of the cumulative emissions that the world could have produced before crossing the 350 ppm boundary. That approach takes historical responsibility into account and makes the boundary much stricter for countries that industrialised early. \citet{hickel2022} applied a similar fair-shares logic to resource use. The equal-per-capita annual boundary used here is therefore a lenient choice for high-income countries. Results that show overshoot under it would only be strengthened under a fair-shares boundary.

The choice of carbon budget also matters. Since \citet{oneill2018}, the IPCC has updated its estimates. The remaining budget from the beginning of 2020 for a 50\% chance of limiting warming to 1.5\textdegree C is about 500 Gt \COtwo{} \citep{ipcc2021spm}. Spread evenly over 2020--2100 and the 2020 population, this implies about \nBoundaryFifteen~\tpc, roughly half the 2\textdegree C value. This paper uses it as a stricter alternative.

\subsection{The social foundation}

The planetary boundaries define what humanity must avoid. They say nothing about what people need. \citet{raworth2012} proposed to complement them with a social foundation, a set of minimum standards drawn from internationally agreed social priorities. The combination of the two, the ``doughnut'', has since become a widely used framework for thinking about sustainable development \citep{raworth2017}. The 2030 Agenda for Sustainable Development \citep{un2015sdg} provides many of the targets used to define the foundation, including the target of reducing under-5 mortality to at most 25 deaths per 1{,}000 live births by 2030.

\citet{oneill2018} defined eleven social thresholds that could be measured across countries. They include life satisfaction, healthy life expectancy, nutrition, sanitation, income above a poverty line, access to energy, secondary education, social support, democratic quality, equality and employment. \citet{fanning2022} used a closely related set of eleven indicators and thresholds, including life expectancy of 74 years, access to electricity for 95\% of the population and employment of 94\% of the labour force. Both studies measured employment as one minus the unemployment rate, using harmonised World Bank data, and set its threshold at 6\% unemployment. The social thresholds in this paper are taken directly from these studies, except for under-5 mortality, which uses the Sustainable Development Goal target.

A related literature asks what material conditions are needed for a decent life. \citet{rao2018} defined ``decent living standards'' in terms of concrete material requirements such as housing, nutrition, health care, education and mobility. \citet{millwardhopkins2020} estimated that decent living standards could be provided for a global population of ten billion with roughly 60\% less final energy than the world uses today, given efficient technologies and equitable provision. \citet{kikstra2021} quantified the gaps in decent living standards across countries and the energy needed to close them. This work suggests that the energy and carbon needed to meet basic needs is, in principle, much lower than what high-income countries currently use. Whether countries have actually achieved basic needs at low emission levels is an empirical question, and one this paper addresses.

\subsection{Human well-being and emissions}

A large body of research documents a strongly nonlinear relationship between energy use or emissions and human well-being. Across countries, life expectancy, education and other basic outcomes rise steeply with energy use at low levels and then level off. At the highest levels of energy use, additional consumption is associated with little or no further gain \citep{steinberger2010}. \citet{steinberger2010} showed that between 1975 and 2005 the energy and carbon associated with given levels of life expectancy and literacy fell substantially. They called this a decoupling of human needs from energy and carbon. \citet{lamb2015} examined what the historical relationship between human development and emissions implies for development within a carbon-constrained world, and how that relationship has shifted over time.

\citet{jorgenson2014} introduced the ``carbon intensity of well-being'', the ratio of emissions per person to life expectancy, and showed that its relationship with economic development differs across regions of the world. \citet{steinberger2020} contrasted the close coupling of emissions with income with the much looser coupling of emissions with life expectancy, and described the resulting tension as a ``carbon-development paradox''. \citet{vogel2021} found that the level of energy use at which human needs are satisfied depends strongly on how societies provide for those needs. Factors such as the quality of public services and access to electricity are associated with higher need satisfaction at lower energy use. Most recently, \citet{porcher2026} described a ``champagne curve'' in which human development rises with consumption-based emissions per person at low levels and then flattens, so that beyond a certain level additional emissions no longer improve well-being.

A complementary tradition in environmental sociology and economics relates emissions to their drivers, such as population and affluence \citep{dietz1997}, and asks whether environmental pressures eventually decline as countries grow richer. This environmental Kuznets curve hypothesis was proposed by \citet{grossman1995}. It has found some support for local pollutants but little for \COtwo{} \citep{stern2004}. This paper does not test the Kuznets hypothesis. The literature is relevant because it shows that the relationship between income and emissions is not fixed, and can change over time and differ across countries.

\subsection{Decoupling emissions from income}

Decoupling refers to the weakening of the link between economic activity and environmental pressure. Relative decoupling occurs when emissions grow more slowly than income. Absolute decoupling occurs when emissions fall while income grows \citep{wiedenhofer2020}. A systematic review of the empirical evidence by \citet{haberl2020} concluded that relative decoupling is common. Absolute decoupling has been observed mainly for greenhouse-gas emissions in some countries and periods, and the observed rates of absolute decoupling are far too low to meet climate targets. \citet{hickel2020green} argued, on the basis of this and related evidence, that green growth is unlikely to be achieved fast enough.

\citet{vogel2023} compared achieved and required rates directly. They identified 11 high-income countries that had absolutely decoupled consumption-based \COtwo{} emissions from GDP between 2013 and 2019. At the rates they achieved, these countries would take on average more than 220 years to reduce their emissions by 95\%, far longer than the Paris Agreement allows. This paper applies a similar comparison to every country with consumption-based emissions data, using emissions per person and the equal-per-capita boundary.

\subsection{Production- and consumption-based accounting}

National emissions can be measured in two ways. Production-based (or territorial) accounting counts the emissions released within a country's borders. Consumption-based accounting adjusts territorial emissions for international trade. It adds the emissions embodied in imports and subtracts those embodied in exports, so that emissions are attributed to the country where the goods and services are finally consumed. \citet{peters2011} showed that emissions embodied in trade grew rapidly between 1990 and 2008. For many high-income countries, the reduction in territorial emissions was partly offset by growing imports from countries with emissions-intensive production.

The choice matters for comparisons with a per-person boundary, all the more because affluent consumption is a major driver of environmental pressure \citep{wiedmann2020}. Production-based accounting credits importing countries with reductions that may simply reflect production moving abroad, and assigns the corresponding emissions to exporters. Consumption-based accounting assigns responsibility to the final users of goods and services, which is closer to the idea of a fair per-person share of the global budget. The Global Carbon Budget publishes consumption-based estimates for most large economies, derived from multi-regional input--output models \citep{friedlingstein2026}. This paper uses consumption-based emissions as its main measure and production-based emissions as a robustness check. Consumption-based estimates are modelled rather than measured and carry additional uncertainty, which is discussed in \cref{sec:discussion}.

\subsection{Gaps addressed by this paper}

The literature reviewed above leaves four gaps that this paper addresses.

\begin{enumerate}
\item \textbf{The dynamics of meeting the social foundation.} National doughnut studies have documented that no country is in the safe and just space at a point in time, and that shortfall and overshoot have both changed since the 1990s. They have not systematically traced the countries that crossed the social foundation and asked on which side of the climate boundary they crossed it.

\item \textbf{The changing carbon cost of meeting basic needs, with uncertainty.} Studies of the energy and carbon intensity of well-being show that the relationship between emissions and well-being has shifted over time. They have rarely summarised this shift in a single quantity tied directly to a set of social thresholds and a climate boundary, with confidence intervals that account for the dependence of observations within countries.

\item \textbf{The adequacy of decoupling beyond high-income countries.} The comparison of achieved and required decoupling rates has focused on high-income countries. Applying it to all countries with consumption-based data shows how widespread the gap is.

\item \textbf{Transparency and robustness.} Results in this field depend on many choices: thresholds, boundaries, accounting methods and the treatment of missing data. This paper reports how each headline result changes under alternative choices, and distinguishes conclusions that are stable from estimates whose level depends on the specification.
\end{enumerate}
